%% ===========================================================================
%% TAILORED CV - Marcus D. Yang - Saronic Manufacturing Engineer Intern
%% ---------------------------------------------------------------------------
%% This file is the FROZEN FORMATTING STANDARD and the FACT SOURCE for every
%% tailored CV. Two rules:
%%   1. The preamble below (geometry, font, \section rule, \erow macro, list
%%      spacing) is frozen. A tailored CV COPIES it verbatim - never retune it.
%%   2. Everything between \begin{document} and \end{document} is CONTENT and
%%      is fully re-orderable / re-groupable per the job description. See
%%      .claude/skills/job-application-assistant/05-cv-templates.md.
%% Compile: lualatex -interaction=nonstopmode main_example.tex   (MUST be 1 page)
%% Visual reference: cv/Marcus Yang Master Resume.pdf
%% ===========================================================================
\documentclass[11pt,letterpaper]{article}

% --- Fonts: Times New Roman clone, bundled with every modern TeX dist -------
\usepackage{fontspec}
\setmainfont{TeX Gyre Termes}
\usepackage{microtype}
\linespread{0.93}   % match the master's tight (Word "single") leading

% --- Page geometry (master: 0.8in sides, ~0.5in top) ----------------------
\usepackage[top=0.5in,bottom=0.5in,left=0.8in,right=0.8in]{geometry}

% --- Section headings: bold, full-width rule underneath --------------------
\usepackage{titlesec}
\titleformat{\section}{\fontsize{12.5}{14}\selectfont\bfseries}{}{0pt}{}[{\vspace{1pt}\hrule height 0.6pt}]
\titlespacing*{\section}{0pt}{6pt}{2pt}

% --- Tight bullet lists ---------------------------------------------------
\usepackage{enumitem}
\setlist[itemize]{leftmargin=1.35em,itemsep=1pt,topsep=1pt,parsep=0pt,
                  partopsep=0pt,label=\textbullet}

\usepackage[hidelinks]{hyperref}
\pagestyle{empty}
\setlength{\parindent}{0pt}

% --- \erow{left}{right}: entry/degree row, left flush-left, right flush-right
%     Caller supplies ALL markup in {left} (\textbf{...}, \textit{...}, etc.).
\newlength{\erowdate}\setlength{\erowdate}{2.0in}
\newcommand{\erow}[2]{%
  \noindent\makebox[\dimexpr\linewidth-\erowdate\relax][l]{#1}%
  \makebox[\erowdate][r]{#2}\par}

% --- \entrygap: vertical space between entries within a section -----------
\newcommand{\entrygap}{\vspace{2pt}}

\begin{document}

% ===========================================================================
%  HEADER
% ===========================================================================
\begin{center}
  {\fontsize{17}{20}\selectfont\bfseries Marcus D. Yang}\\[2pt]
  {\footnotesize (650) 954-0288 \textbar{} \href{mailto:yangdmarcus@gmail.com}{yangdmarcus@gmail.com} \textbar{} \href{https://www.linkedin.com/in/marcusdyang}{www.linkedin.com/in/marcusdyang} \textbar{} \href{https://projects-website-psi.vercel.app/}{https://projects-website-psi.vercel.app/}}
\end{center}

% ===========================================================================
%  EDUCATION
% ===========================================================================
\section{Education}
\erow{\textbf{University of Illinois Urbana-Champaign} \textbar{} Champaign, IL}{August 2024 - May 2028}
\erow{\textit{B.S. in Aerospace Engineering with a Computer Science minor, James Scholar Honors}}{GPA: 3.57/4}
\textbf{Relevant Coursework:} Engineering Materials, Aerospace Structures, Statics, Dynamics, Thermodynamics, Electrical and Electronic Circuits, Aerospace Control Systems, Incompressible Flow, Linear Algebra, Discrete Structures

% ===========================================================================
%  ORGANIZATIONAL AND TECHNICAL EXPERIENCE
% ===========================================================================
\section{Organizational and Technical Experience}

\erow{\textbf{TVC Rocket} (Fabrication, 3D Printing)}{Summer 2026}
\begin{itemize}
  \item Designed and fabricated a 2-axis TVC gimbal, avionics bay, and recovery system, integrating a Teensy 4.1 flight computer, MPU-6050 IMU, and servo actuators into a commercial rocket body.
  \item Designed engine mount and structural parts in Siemens NX (NXOpen/Python) for FDM additive manufacturing.
  \item Performed root cause analysis on two static fire test failures from telemetry and video data, isolating a misreferenced attitude axis and an inverted control loop sign.
  \item Traced a servo linkage resolution bottleneck by characterizing the mechanical actuation ratio, then redesigned the control path to reach sub-degree pointing precision.
  \item Bench tested the PID control loop against a live IMU feed to confirm gimbal response before static fire.
  \item Logged telemetry at 83 Hz with zero dropouts across two static fires (15,728 and 12,611 samples).
  \item Ran a Monte Carlo sensitivity analysis across hundreds of trials (IMU noise, motor burn scatter, crosswind), finding landing performance was more sensitive to accelerometer bias than servo speed and redirecting effort toward sensor calibration.
\end{itemize}
\entrygap

\erow{\textbf{Liquid Rocketry at Illinois} (Siemens NX, Assembly Integration)}{Spring 2025 - Present}
\begin{itemize}
  \item Modeled P\&ID system architecture in Siemens NX and integrated components into the pressurant tank assembly.
  \item Assessed pressurant tank options and produced mechanical integration models verifying structural fit.
  \item Engineered a KiCad PCB radio transmitter (USB-C + MCU) prototype enabling ground-test telemetry links.
  \item Created CAD geometries for canard-grid fins and ran Ansys Fluent CFD cases across 0--20\textdegree{} deflection.
\end{itemize}
\entrygap

\erow{\textbf{Model Rocket Project} (Siemens NX, OpenRocket)}{Fall 2024}
\begin{itemize}
  \item Designed custom fin geometry in Siemens NX and fabricated a physical prototype to hit a target apogee.
  \item Validated against OpenRocket simulations; flight reached 430 ft apogee with 25 ft downrange drift.
\end{itemize}
\entrygap

\erow{\textbf{Glider Project}}{Spring 2026}
\begin{itemize}
  \item Fabricated and iteratively tested gliders, tuning mass distribution and wing geometry for lift-to-drag performance.
  \item Achieved 101 ft glide distance and 3.9 s flight time, ranking 6th in class.
\end{itemize}

% ===========================================================================
%  TECHNICAL SKILLS
% ===========================================================================
\section{Technical Skills}
\textbf{Computer Languages:} Python, C++, Java, JavaScript\\
\textbf{Software:} Siemens NX, Fusion 360, Eagle PCB, KiCad, OpenRocket, XFLR5, Ansys Fluent, Excel\\
\textbf{Shop Tools:} Bambu A1 Mini, Spot Welder, Belt Sander, Laser Cutter, CNC Plasma Cutter, Table Saw, Drill Press\\
\textbf{Materials:} Epoxy Resin, Fiberglass, Sheet Metal, Aluminum Pipes, Plywood, PETG, PLA

% ===========================================================================
%  LEADERSHIP AND ACTIVITIES
% ===========================================================================
\section{Leadership and Activities}

\erow{\textbf{Flight Club Aerospace}}{Summer 2021 - Spring 2024}
Landing Gear Research Lead \textbar{} CAD \textbar{} Fabrication
\begin{itemize}
  \item Manufactured wing and control surfaces: cut foam ribs, trimmed aluminum spars with a circular saw, and applied epoxy \& fiberglass layups to shape aerodynamic skins; aligned two 7-ft wings for flight testing.
  \item Drafted gussets and ribs in Fusion 360 and confirmed fit-checks between structural and aerodynamic components.
  \item Oversaw logistics for wing transport/disassembly/reassembly and assembled and presented landing-gear recommendations based on structural load and cost analyses.
\end{itemize}
\entrygap

\erow{\textbf{Asian American Association}}{Fall 2024 - Present}
Treasurer \textbar{} Fundraising Officer
\begin{itemize}
  \item Managed membership operations for 300+ members, coordinating records and organizational logistics.
  \item Secured partnerships with 15 on-campus sponsors to support events and fundraising initiatives.
\end{itemize}

\end{document}
